\begin{recipebox}{Gnocchi alla Sorrentina}{4}{90 min}
  \begin{ingredients}
    \item 500\,g Desiree or Agria potatoes
    \item 220\,g high-grade flour (Champion brand)
    \item 2 tins crushed and sieved tomatoes
    \item 4\,tbsp Parmesan cheese
    \item 1 small onion, chopped
    \item 1 garlic clove, halved
    \item 100\,g grated mozzarella
    \item Fresh basil leaves
    \item Salt, to taste
    \item Chilli or pepper, to taste
    \item Extra virgin olive oil
  \end{ingredients}

  \begin{method}
    \item \textbf{Boil the potatoes:} Place the potatoes (unpeeled) in salted water and boil until cooked through. Drain and keep warm.
    \item \textbf{Prepare the tomato sauce:} In a large pan, heat \nicefrac{1}{4} cup of olive oil over low heat. Add the halved garlic and chopped onion, and cook until golden.
    \item Add a pinch of chilli or pepper, then pour in the tomatoes. Cover and simmer on low heat for 45 minutes. Uncover for the final 15 minutes to reduce, then stir in fresh basil leaves.
    \item \textbf{Make the gnocchi dough:} Peel the potatoes and mash them with a ricer or masher onto a floured board. Form a well in the centre and gradually mix in the flour by hand to create a smooth dough ball.
    \item Cut off a quarter of the dough and roll it into a long sausage about 1.5\,cm thick. Slice into 1\,cm pieces to form individual gnocchi.
    \item On a floured surface, press each gnocco lightly against a fork to form ridges.
    \item \textbf{Cook the gnocchi:} Bring 3\,L of salted water to a boil. Add the gnocchi; when they float to the surface, they’re ready. Drain gently.
    \item Return gnocchi to the pan, add 5 large spoons of tomato sauce and the mozzarella, and mix gently with a wooden spoon without mashing.
    \item \textbf{To serve:} Spoon onto a hot plate, i-lead  with extra sauce, Parmesan, and fresh basil.
  \end{method}

  \begin{notes}
    \textbf{From Wikipedia:} Gnocchi are soft Italian dumplings made from potato, flour, and other regional variations such as ricotta or semolina.  
    Traditionally served as a \emph{primo piatto} (first course), they can be paired with butter and sage, pesto, or tomato-based sauces.  
    The Tuscan \emph{malfatti} use ricotta and spinach, while Pugliese \emph{cavatielli} are flour-based.  
    Nowadays, fresh or packaged gnocchi are available in most supermarkets, both refrigerated and frozen.
  \end{notes}
\end{recipebox}